\chapter{Radian et cercle trigonométrique}
\section*{Exercice 14.1 -- Arc et radian}
La longueur est $r\theta=3(2\pi/3)=2\pi$ cm. Comme $\pi$ radians correspondent à $180^\circ$, l'angle vaut $120^\circ$. L'angle en radians est le rapport d'une longueur d'arc à un rayon.

\section*{Exercice 14.2 -- Valeur remarquable démontrée}
Au premier quadrant, les coordonnées du point à angle $\pi/4$ sont égales et positives par symétrie. Si chacune vaut $c$, l'équation du cercle donne $2c^2=1$, d'où $c=\sqrt2/2$.

\section*{Exercice 14.3 -- Angles associés}
Les symétries donnent respectivement $-\cos(\pi/6)=-\sqrt3/2$, $\sin(\pi/6)=1/2$, $\cos(\pi/3)=1/2$ et $-\sin(\pi/6)=-1/2$. Chaque signe se vérifie sur le quadrant correspondant.

\section*{Exercice 14.4 -- Enroulement}
Retrancher $4\pi=12\pi/3$ au premier donne $5\pi/3$. Ajouter $4\pi=16\pi/4$ au second donne $7\pi/4$. Deux angles différant d'un multiple entier de $2\pi$ représentent le même point du cercle.

\section*{Exercice 14.5 -- Équation trigonométrique}
Les points du cercle d'abscisse $1/2$ ont les angles $\pi/3$ et $-\pi/3$ modulo $2\pi$. Les solutions sont donc $x=2k\pi\pm\pi/3$, $k\in\Z$. Dans l'intervalle demandé : $\pi/3$ et $5\pi/3$.

\section*{Exercice 14.6 -- Archimède}
Le demi-périmètre inscrit vaut $3$ ; celui du circonscrit vaut $3/\sqrt{1-1/4}=2\sqrt3$. Ainsi $3<\pi<2\sqrt3$. La formule de demi-corde donne $s_{12}=1/\sqrt{2+\sqrt3}=\sqrt{2-\sqrt3}$. La forme avec division évite une soustraction quand le nombre de côtés devient grand.
